\section{A logarithmic endpoint on the original source}
\label{sec:v89-logarithmic}
Keep $\gamma=1/32$ and $Q=32/31$. Endpoint boundedness can be
combined with small mass without exchanging a microscopic supremum
and a collision limit. The resulting estimate acts on the original
critical-or-short-residence union, not only on its complement.

\begin{proposition}[Uniform critical weak bound for the raw density]
\label{prop:v89-density-bound}
For every fixed $L>0$,
\begin{equation}\label{eq:v89-density-bound}
 \limsup_m\sup_{R,n,k}\sup_{|J|=L}
                      \mathcal N_Q(m^2p_{n,R}(k,m,\cdot);J)\le K_L.
\end{equation}
For the raw error $f_{m,R,n,k}=m^2p-\mathcal L$, the same assertion
holds with $f_{m,R,n,k}\one_{\mathcal C(M)}$ and a constant $K_{M,L}$.
\end{proposition}
\begin{proof}
Use the source comparison at auxiliary guard parameter one. Its
first factor is zero since $c_0\le1$, so $W^1=1$ and $b^1=p$.
Equations~\eqref{eq:v83-shift-input} and
\eqref{eq:v88-optical-input} give
\[
 p\le(2+C\sqrt r)A_rp+C\overline{\mathcal F}_r
                         +C\mathsf P_{1/4,r}v^{1/4}.
\]
Here $r>0$ is fixed and small; every source on the right is complete.
The fixed-width input bounds the first term in essential height.
Corollary~\ref{cor:v89-full-critical} bounds the second in
$\mathcal N_Q$. Minkowski gives, since $\kappa Q<1$,
\[
 \|\mathsf P_{1/4,r}v\|_{L^Q(J)}
 \le (1-\kappa Q)^{-1/Q}r^{1/Q-\kappa}
                                      \int_{J+(0,r)}v.
\]
H\"older implies $\mathcal N_Q\le\|\cdot\|_{L^Q}$.
The actual fractional fixed-width input therefore bounds the last
term. Take the collision limsup at fixed widths and auxiliary bands,
and enlarge the bands afterwards. This proves the first assertion.
For the error, use the positive scalar remainder at one fixed
sufficiently large reconstruction band. Its positive term is bounded
by $m^2p$, and its other term has bounded central essential height.
No pointwise estimate for $p$ or division by the arithmetic reference
is used.
\end{proof}

For $|J|=L$ and $\sigma>0$ define the logarithmically weakened
critical measurable-set norm
\begin{equation}\label{eq:v89-log-norm}
 \mathcal N_{Q,\sigma}(f;J)
 =\sup_{E\subset J,\ 0<|E|\le L}
 \frac{|E|^{-\gamma}\int_E|f|}
                       {(1+\log(L/|E|))^\sigma}.
\end{equation}
The logarithmic denominator is part of the conclusion, not a
change to the density, event, or physical guard.

\begin{lemma}[Endpoint interpolation with the mass]
\label{lem:v89-log-interpolation}
Suppose $\mathcal N_Q(f;J)\le A$, $A>0$, and put
$t=\|f\|_{L^1(J)}/(AL^\gamma)$. Then $0\le t\le1$.
For $0<t<1$,
\begin{equation}\label{eq:v89-log-interpolation}
 \mathcal N_{Q,\sigma}(f;J)
 \le A\left\{t^{1/2}
       +\left(1+\frac{\log(1/t)}{2\gamma}\right)^{-\sigma}\right\}.
\end{equation}
If $t=0$, then $f=0$ almost everywhere and the left side is zero.
\end{lemma}
\begin{proof}
Split the tested sets at measure $\delta=Lt^{1/(2\gamma)}$.
For $|E|\ge\delta$, their numerator is at most
$\delta^{-\gamma}\|f\|_1=At^{1/2}$. For smaller sets use the
bound $A$ for the numerator and
$1+\log(L/|E|)\ge1+\log(1/t)/(2\gamma)$ in the denominator.
The maximum of the two bounds is at most their sum. This is an
inequality for each actual function before taking any limit.
\end{proof}

\begin{theorem}[The raw error at the logarithmic critical endpoint]
\label{thm:v89-log-raw}
For every fixed $M,L,\sigma>0$,
\begin{equation}\label{eq:v89-log-raw}
 \sup_{R,n,k}\sup_{|J|=L}
 \mathcal N_{Q,\sigma}(f_{m,R,n,k}\one_{\mathcal C(M)};J)
                                                   \longrightarrow0.
\end{equation}
In particular the critical exponent is attained with any positive
logarithmic loss. Convergence in $\mathcal N_Q$ without that loss
is not asserted.
\end{theorem}
\begin{proof}
The subcritical raw theorem includes uniform local $L^1$
convergence. Proposition~\ref{prop:v89-density-bound} gives an
eventually uniform bound $A=2K_{M,L}+1$ for the critical norm.
Apply Lemma~\ref{lem:v89-log-interpolation} to each actual error,
then take the parameter and interval suprema. Their mass ratios
tend uniformly to zero. Equivalently the inherited remainder gives
an ordered $L^1$ bound $C_{M,L}B^{-1/192}$ for each fixed outer
band; inserting it in the lemma and sending $B$ to infinity last
gives the same conclusion. No $B=B(m)$ is chosen.
\end{proof}

\subsection{The exceptional union itself}
Set $a_h=h^{1/8}$, $s_h=h^{1/4}$ and define the original weighted
restriction
\begin{equation}\label{eq:v89-residual-source}
 R_h^\varepsilon(t)
  =\int W^\varepsilon(x)
       \one_{\{z(x)<a_h\ \mathrm{or}\ \ell(x)<s_h\}}\,d\mu_t(x).
\end{equation}
The two conditions form one union in this definition. All endpoint
types and all original exact labels remain in the source.

\begin{theorem}[A critical logarithmic estimate on the exceptional source]
\label{thm:v89-exceptional-log}
For every fixed $L,\sigma>0$, uniformly in $0<\varepsilon\le1$,
\begin{align}
 \limsup_m\sup_{R,n,k}\sup_{|J|=L}
        m^2\int_J R_h^\varepsilon&\le C_Lh^{1/128},
                                    \label{eq:v89-residual-mass}\\
 \limsup_m\sup_{R,n,k}\sup_{|J|=L}
   \mathcal N_{Q,\sigma}(m^2R_h^\varepsilon;J)
                   &\le C_{L,\sigma}(1+|\log h|)^{-\sigma}.
                                    \label{eq:v89-residual-log}
\end{align}
All suprema over tested measurable sets are part of the norm before
the collision limsup. The second estimate is not an essential-height
bound and does not prove the residual charge zero.
\end{theorem}
\begin{proof}
On $0<z<a$, one has
$1\le a^{1/4}I_m z^{-1/4}$. The critical fractional input therefore
gives ordered interval mass at most $C_La^\alpha$, with
$\alpha=129/128$. The complete short-residence bound gives at most
$C_Ls^\gamma$. Since $W^\varepsilon\le1$, a positive union bound
at $a=a_h$, $s=s_h$ proves
\[
 C_L(h^{129/1024}+h^{1/128})\le C_L' h^{1/128}.
\]
This upper comparison does not double count the source definition.
Also $R_h^\varepsilon\le p$, so its critical weak norm is bounded
by Proposition~\ref{prop:v89-density-bound}, independently of $h$
and $\varepsilon$. Fix $h$ before the collision limsup and apply
Lemma~\ref{lem:v89-log-interpolation}. The mass ratio is at most
$C_Lh^{1/128}$ in that limit. Its power term is absorbed by the
logarithmic right side in \eqref{eq:v89-residual-log}.
The estimate applies to the exceptional union itself, including
critical initial points in every positive neighborhood and arbitrarily
short regular segments; it uses no common-survivor restriction.
\end{proof}

The cofinal limit of the essential heights of
\eqref{eq:v89-residual-source} is the same original charge as in
the retained source-localization theorem: both positive thresholds
decrease to zero, and the restriction is monotone in each threshold.
Thus the governing pointwise assertion still requires that charge to
vanish in the outer reconstruction limit. A logarithmic measurable-set
estimate is an additional quantified restriction on concentration,
not a proof that its essential height vanishes.
